% Waterhose coupled-env scaling report -- single-file Overleaf template.
% Data: _scratch/reports/waterhose_scaling/waterhose_env_scaling_2026-06-24.csv (scaling sweep, all
%       collisions on) + profile_components.py stdout (single-env per-step breakdown).
% Styling adapted from _scratch/reports/scoop_env_scaling_report.tex (IsaacLab-coupling).
\documentclass[a4paper,11pt]{article}

\usepackage[margin=0.85in, top=1.3in, bottom=0.85in, headheight=18pt, headsep=14pt]{geometry}
\usepackage[T1]{fontenc}
\usepackage[utf8]{inputenc}
\usepackage[scaled=0.95]{helvet}
\usepackage{microtype}
\usepackage{xcolor}
\usepackage{booktabs}
\usepackage{array}
\usepackage{tabularx}
\usepackage{graphicx}
\usepackage{siunitx}
\sisetup{
  group-separator={,},
  group-minimum-digits=4,
  detect-weight=true,
  detect-family=true,
  table-number-alignment=center,
}
\usepackage{enumitem}
\usepackage{amsmath}
\usepackage{amssymb}
\usepackage{pgfplots}
\usepackage{tikz}
\usepackage{tcolorbox}
\tcbuselibrary{skins,breakable}
\usepackage{titlesec}
\usepackage{fancyhdr}
\usepackage{hyperref}

\renewcommand{\familydefault}{\sfdefault}
\pgfplotsset{compat=1.18}
\usetikzlibrary{calc}

% NVIDIA / report palette (matches Isaac Lab ecosystem)
\definecolor{NvidiaGreen}{HTML}{76B900}
\definecolor{NvidiaGreenDark}{HTML}{558B00}
\definecolor{ReportText}{HTML}{1A1A1A}
\definecolor{ReportMuted}{HTML}{4D4D4D}
\definecolor{SummaryGrey}{HTML}{F2F2F2}
\definecolor{RuleGrey}{HTML}{D9D9D9}
\definecolor{TabBlue}{HTML}{1F77B4}
\definecolor{TabGreen}{HTML}{2CA02C}
\definecolor{TabOrange}{HTML}{FF7F0E}
\definecolor{TabRed}{HTML}{D62728}

\color{ReportText}
\hypersetup{colorlinks=true, linkcolor=NvidiaGreenDark, urlcolor=NvidiaGreenDark}

\setlength{\parindent}{0pt}
\setlength{\parskip}{0.45em}
\setlength{\headheight}{22pt}
\setlength{\headsep}{18pt}

\newcommand{\NvidiaTopBanner}{%
  \tikz[remember picture,overlay]{%
    \fill[NvidiaGreen] (current page.north west) rectangle ([yshift=-1.05cm]current page.north east);
    \node[anchor=west, text=white, font=\bfseries\small]
      at ([xshift=0.85in, yshift=-0.55cm]current page.north west) {NVIDIA | ISAAC LAB};
    \node[anchor=east, text=white, font=\small]
      at ([xshift=-0.85in, yshift=-0.55cm]current page.north east) {Internal technical report};
  }%
}

\newcommand{\ReportRunHeader}{%
  \begin{tcolorbox}[
    enhanced, colback=NvidiaGreen, colframe=NvidiaGreen, boxrule=0pt, arc=0pt,
    left=6pt, right=6pt, top=4pt, bottom=4pt, width=\headwidth, halign=left,
  ]
  {\color{white}\footnotesize\textbf{NVIDIA | ISAAC LAB}\hfill
   Scaling Report | Waterhose Coupled (RBY1DF)}
  \end{tcolorbox}%
}

\newcommand{\ReportRule}{{\color{RuleGrey}\rule{\linewidth}{0.5pt}}}
\newcommand{\tableheading}[1]{{\color{ReportMuted}#1\par\vspace{0.35em}}}
\newcommand{\ReportCaption}[1]{{\small\color{ReportMuted} #1\par}}

\newenvironment{reporttable}{%
  \vspace{0.35em}\begin{center}\small
  \renewcommand{\arraystretch}{1.3}\setlength{\tabcolsep}{10pt}%
}{%
  \end{center}\vspace{0.55em}%
}

\newcommand{\ReportHeading}[1]{%
  \vspace{0.6em}{\color{NvidiaGreen}\bfseries\large #1\par}\vspace{0.25em}\ReportRule\vspace{0.45em}%
}

\newtcolorbox{executivesummarybox}{
  enhanced, breakable, colback=SummaryGrey, colframe=SummaryGrey, boxrule=0pt,
  borderline west={3.5pt}{0pt}{NvidiaGreen}, left=10pt, right=10pt, top=6pt, bottom=6pt,
  arc=0pt, before skip=0.3em, after skip=0.55em,
}

\titleformat{\section}
  {\color{NvidiaGreen}\bfseries\large}{\thesection\ }{0pt}{}
  [\vspace{0.2em}\ReportRule\vspace{0.35em}]
\titlespacing*{\section}{0pt}{0.9em}{0.5em}

\pagestyle{fancy}
\fancyhf{}
\fancyhead[C]{\ReportRunHeader}
\fancyfoot[C]{\small\thepage}
\renewcommand{\headrulewidth}{0pt}
\fancypagestyle{firstpage}{\fancyhf{}\fancyfoot[C]{\small\thepage}\renewcommand{\headrulewidth}{0pt}}

\pgfplotsset{
  reportplot/.style={
    width=0.86\linewidth, height=6.2cm,
    tick label style={/pgf/number format/1000 sep={}},
    label style={font=\small}, title style={font=\normalsize\bfseries},
    grid=both, grid style={gray!25}, major grid style={gray!25},
    minor tick num=1, minor grid style={gray!12},
    legend style={font=\small, fill=white, draw=gray!40},
    enlarge x limits=0.03, enlarge y limits=0.06,
  },
  wh envaxis/.style={
    xmode=log, log basis x={2}, xtick pos=left, xlabel style={font=\small},
    xmin=440, xmax=9400, xtick={512,1024,2048,4096,8192}, xticklabels={512,1024,2048,4096,8192},
  },
}

\newcommand{\B}{\bfseries}

\begin{document}
\thispagestyle{firstpage}
\NvidiaTopBanner

\vspace*{0.05cm}
{\small\color{ReportMuted}%
June 24, 2026 \quad|\quad Branch \texttt{max/waterhose-coupled-experimental}
\quad|\quad Prepared by Maximilian Krause\par}

\vspace{0.55em}
{\fontsize{20}{24}\selectfont\bfseries Waterhose Coupled Environment Throughput Scaling\par}
\vspace{0.2em}
{\normalsize\color{ReportMuted} \texttt{Isaac-Waterhose-Coupled-v0} --- proxy vs.\ ADMM coupling,
all collisions active, 512--8192 worlds\par}

\vspace{0.4em}\ReportRule\vspace{0.45em}

\ReportHeading{Executive summary}

\begin{executivesummarybox}
\begin{itemize}[leftmargin=1.25em, itemsep=0.15em, topsep=0pt, parsep=0pt]
  \item \textbf{Proxy is the faster backend at every world count, and scales to 6144} ---
        \textbf{32{,}544 env-steps/s at 6144 worlds} (the peak) on a 32\,GiB RTX~5090, all CUDA-graph
        captured. Proxy climbs monotonically (10.9k$\to$17.9k$\to$24.7k$\to$29.2k$\to$31.2k$\to$32.5k
        from 512 to 6144); 8192 crashes (NaN at $\approx$\,26\,GiB --- the deterministic-contact cap,
        not OOM).
  \item \textbf{ADMM peaks at 1024 ($\approx$\,12.3k env-steps/s) and is $\approx$\,2.6$\times$ slower
        than proxy at scale.} Proxy resolves contacts through Newton's collision pipeline; ADMM disables
        it (\texttt{use\_collision\_pipeline=False}). ADMM is also heavier --- 15.5\,GiB at 2048 vs
        proxy's 12.2.
  \item \textbf{Memory tops out near the card, startup is steep.} Proxy reaches 27.5\,GiB at 6144 (near
        the 32\,GiB ceiling); the hose is a thin Cosserat rod so per-env memory is small, but
        \emph{build/startup} grows super-linearly ($\approx$\,24\,s at 512 to $\approx$\,14\,min at 6144).
        Per step the VBD/AVBD cable+contact solve dominates (40\,\%); the CUDA graph is essential
        ($\approx$\,6$\times$).
\end{itemize}
\end{executivesummarybox}

\section{Methodology}

Each point is one headless, kitless, CUDA-graph run on an idle NVIDIA GeForce RTX~5090, 40 warm-up + 60
measured steps, driven by the scripted grasp-insert policy (eight coupled substeps per step: MuJoCo-Warp
arm, VBD cable + hard AVBD contacts, coupling, plus the IK solve). The collision config is active for
both backends: robot$\leftrightarrow$housing (a single welded concave mesh, collided through Newton with
\texttt{use\_mujoco\_contacts=False}), and plug$\leftrightarrow$socket. \textbf{Proxy}
(\texttt{Isaac-Waterhose-Coupled-v0})
mirrors the gripper as proxy bodies and resolves contacts through Newton's collision pipeline;
\textbf{ADMM} (\texttt{Isaac-Waterhose-Admm-v0}) couples through an ADMM contact constraint with the
pipeline disabled. ADMM was measured to 2048 here; proxy ran the full 512--8192 range. 8192 crashes for
proxy (NaN, not memory).

\section{Scaling results}

\tableheading{Throughput and peak GPU memory vs.\ world count (all collisions on, CUDA graph on).
Startup is the scene build (\texttt{gym.make}); $\approx$ for $\ge$\,3072 (derived from wall time).}
\begin{reporttable}
\setlength{\tabcolsep}{8pt}%
\begin{tabular}{@{}r cc cc c@{}}
\toprule
& \multicolumn{2}{c}{\textbf{Proxy (default)}} & \multicolumn{2}{c}{\textbf{ADMM}} & \\
\cmidrule(lr){2-3}\cmidrule(lr){4-5}
\textbf{worlds} & {env-steps/s} & {peak GiB} & {env-steps/s} & {peak GiB} & \textbf{startup (s)} \\
\midrule
512  & 10932  & 6.6  & 6726  & 6.9  & 24 \\
1024 & 17706  & 8.5  & 12270 & 8.9  & 47 \\
2048 & 24726  & 12.2 & 11330 & 15.5 & 107 \\
3072 & 29226  & 15.8 & ---   & ---  & $\approx$250 \\
4096 & 31163  & 20.6 & ---   & ---  & $\approx$415 \\
6144 & \B 32544 & 27.5 & ---  & ---  & $\approx$835 \\
8192 & \multicolumn{2}{c}{\emph{crash (NaN @ 26\,GiB)}} & --- & --- & --- \\
\bottomrule
\end{tabular}
\end{reporttable}
\ReportCaption{Proxy peaks at 6144 worlds (32{,}544 env-steps/s, near the 32\,GiB ceiling); 8192 crashes
with a NaN at $\approx$\,26\,GiB (deterministic-contact-matching cap, not OOM). ADMM peaks at 1024 and
was measured only to 2048 (and is heavier --- 15.5\,GiB there). Build/startup grows super-linearly:
$\approx$\,24\,s at 512 to $\approx$\,14\,min at 6144.}

\begin{center}
\begin{tikzpicture}
\begin{axis}[
  reportplot, wh envaxis,
  title={Total throughput vs.\ world count (all collisions on)},
  xlabel={Number of parallel worlds (log scale)}, ylabel={Env-steps per second},
  ymin=0, ymax=35000, ytick={0,5000,10000,15000,20000,25000,30000}, legend pos=north west,
]
\addplot[TabBlue, thick, dashed, mark=o, mark size=2.0pt] coordinates {
  (512,6726) (1024,12270) (2048,11330)
};
\addlegendentry{ADMM coupling}
\addplot[TabGreen, thick, mark=*, mark size=1.8pt] coordinates {
  (512,10932) (1024,17706) (2048,24726) (3072,29226) (4096,31163) (6144,32544)
};
\addlegendentry{Proxy (default)}
\addplot[only marks, TabRed, mark=*, mark size=2.6pt] coordinates {(6144,32544)};
\node[font=\scriptsize, color=TabRed, anchor=north east] at (axis cs:8000,31500) {proxy peak 6144};
\end{axis}
\end{tikzpicture}

\ReportCaption{Proxy rises monotonically and peaks at 6144 worlds (32{,}544 env-steps/s); 8192 crashes
(NaN, not OOM). ADMM peaks at 1024 and was measured only to 2048 --- proxy is faster at every count.}
\end{center}

\section{Optimal world count}

Proxy peaks at \textbf{6144 worlds}; 8192 is not viable (NaN below the memory ceiling). For day-to-day
runs, \textbf{2048--4096} is the better trade-off: most of the peak throughput at a fraction of the build
time (see scaling table above). Use ADMM only when proxy coupling is unavailable --- it is slower at
every measured count.

\section{Per-step cost (single env)}

\tableheading{Where one env-step's GPU time goes, by component (single env, CUDA graph disabled to time
kernels by name; 8 substeps $\times$ 16 VBD iters --- the tuned default). Total kernel time 24.5\,ms/step.}
\begin{reporttable}
\begin{tabular}{@{}l S[table-format=2.1] S[table-format=2.1]@{}}
\toprule
\textbf{Component} & {\textbf{ms/step}} & {\textbf{\% GPU}} \\
\midrule
VBD / AVBD solve (cable + hard contacts) &  9.8 & 40.0 \\
MJWarp (rigid robot)                     &  7.6 & 30.9 \\
Overhead (memset/memcpy)                 &  2.9 & 12.0 \\
Newton IK                                &  2.5 & 10.2 \\
Collision (broad/narrow phase)           &  1.4 &  5.8 \\
Coupling (proxy)                         &  0.3 &  1.1 \\
\bottomrule
\end{tabular}
\end{reporttable}
\ReportCaption{VBD dominates. The MJWarp row is inflated by coupled-solver state-mapping copies and the
Newton-contact path --- the bare articulation is only $\approx$\,0.29\,ms/substep (a standalone
MuJoCo-Warp robot benchmark), so the rigid solve is not the bottleneck. With the CUDA graph enabled the
per-env step is $\approx$\,27\,ms ($\approx$\,kernel sum); graph-disabled the wall step is
$\approx$\,167\,ms, so the graph is worth $\approx$\,6$\times$.}

\section{Recommended operating points}

\begin{reporttable}
\setlength{\tabcolsep}{6pt}%
\begin{tabularx}{\linewidth}{@{}>{\raggedright\arraybackslash}p{0.20\linewidth}
  >{\raggedright\arraybackslash}p{0.30\linewidth} X@{}}
\toprule
\textbf{Goal} & \textbf{Config} & \textbf{Notes} \\
\midrule
Max throughput   & Proxy @ 6144 worlds & Peak env-steps/s; long startup. \\
Fast iteration   & Proxy @ 2048--4096  & Best throughput vs.\ build-time trade-off. \\
Demo / teleop    & Proxy @ 1 env       & Tuned grasp-insert policy. \\
Cheaper step     & \texttt{SUBSTEPS=6}, \texttt{VBD\_ITERS=12} & $\approx$1.44$\times$ per-step speedup vs.\ 8/16 default. \\
\bottomrule
\end{tabularx}
\end{reporttable}

\vspace{0.6em}
{\small\color{ReportMuted}%
Data (2026-06-25, single welded-housing collider, 8/16 substeps): proxy 512--8192, ADMM 512--2048;
drivers \texttt{bench\_sweep.py} / \texttt{bench\_child.py}; per-step profile
\texttt{profile\_components.py}. ADMM $\ge$\,3072 not yet measured.}

\end{document}
